%% kapa: exam=1 round=26 year=2015 period=1 form=A subject=civil_law subject_ko=민법 seq=34 number=34
%% source: Q-Net 감정평가사 기출문제 (cst00309) artlSeq=5138001 file=extracted/1008/5138001/제26회_감정평가사 1차_기출/제26회_감정평가사 1차_기출/1교시_감평_A형.pdf
%% conversion: kordoc 4.24.51 -> tools/gen1.py (auto); review=proofread-sonnet; flags=none
\begin{question}{34}
2014년 丙 소유 X토지를 취득하고 싶은 甲은 그 친구 乙과 X토지의 취득에 관한 명의신탁약정을 맺고 乙에게 X토지를 매수하기 위한 자금을 제공하면서 乙 명의로 丙과 계약하도록 하였다. 이에 乙은 그 사실을 알지 못하는 丙으로부터 X토지를 매수하여 자기 앞으로 이전등기를 마쳤다. 다음 설명으로 옳지 않은 것은? (다툼이 있으면 판례에 따름)
\begin{choices}
  \item 甲과 乙의 명의신탁약정은 무효이다.
  \item 乙은 甲으로부터 받은 X토지 매수대금을 甲에게 부당이득으로 반환할 의무가 있다.
  \item 만약 丙이 명의신탁약정의 존재를 알았다면 X토지에 관한 물권변동은 무효이다.
  \item 만약 乙이 완전한 소유권을 취득했음을 전제로 사후적으로 甲과 매수자금반환의무의 이행에 갈음하여 X토지를 양도하기로 약정하고 甲 앞으로 소유권이전등기를 마쳤다면, 그 등기는 원칙적으로 유효하다.
  \item 만약 甲과 乙의 명의신탁약정 및 乙 명의의 등기가 「부동산 실권리자 명의등기에 관한 법률」의 시행 전에 이루어지고, 같은 법 소정의 유예기간 내에 甲 앞으로 등기가 되지 않았다면, 乙은 甲에게 X토지를 부당이득으로 반환할 의무가 없다.
\end{choices}
\end{question}
\answer{5}
